\documentclass[11pt]{article}
% ===================================================================
% Shared preamble for the "tiny lemma, read slowly" preprint series.
% Mirrors the style of FrobeniusLemmas_preprint.tex.
% ===================================================================

\usepackage{amsmath,amssymb,amsthm}
\usepackage{mathtools}
\usepackage[margin=1.2in]{geometry}
\usepackage{hyperref}
\usepackage{xcolor}
\usepackage{listings}
\usepackage{tikz}
\usepackage{tikz-cd}
\usepackage{mathpartir}
\usepackage{newunicodechar}
\usepackage{setspace}

\onehalfspacing

% ---- Unicode glyphs ----
\newunicodechar{→}{\ensuremath{\to}}
\newunicodechar{↦}{\ensuremath{\mapsto}}
\newunicodechar{↔}{\ensuremath{\leftrightarrow}}
\newunicodechar{⟹}{\ensuremath{\Longrightarrow}}
\newunicodechar{⟶}{\ensuremath{\longrightarrow}}
\newunicodechar{≠}{\ensuremath{\neq}}
\newunicodechar{≤}{\ensuremath{\leq}}
\newunicodechar{≥}{\ensuremath{\geq}}
\newunicodechar{∀}{\ensuremath{\forall}}
\newunicodechar{∃}{\ensuremath{\exists}}
\newunicodechar{∘}{\ensuremath{\circ}}
\newunicodechar{·}{\ensuremath{\cdot}}
\newunicodechar{∈}{\ensuremath{\in}}
\newunicodechar{∉}{\ensuremath{\notin}}
\newunicodechar{≃}{\ensuremath{\simeq}}
\newunicodechar{≅}{\ensuremath{\cong}}
\newunicodechar{≈}{\ensuremath{\approx}}
\newunicodechar{∧}{\ensuremath{\wedge}}
\newunicodechar{∨}{\ensuremath{\vee}}
\newunicodechar{¬}{\ensuremath{\neg}}
\newunicodechar{⊢}{\ensuremath{\vdash}}
\newunicodechar{⟨}{\ensuremath{\langle}}
\newunicodechar{⟩}{\ensuremath{\rangle}}
\newunicodechar{⇑}{\ensuremath{\mathord{\uparrow}}}
\newunicodechar{↑}{\ensuremath{\mathord{\uparrow}}}
\newunicodechar{∎}{\ensuremath{\blacksquare}}
\newunicodechar{∣}{\ensuremath{\mid}}
\newunicodechar{∤}{\ensuremath{\nmid}}
\newunicodechar{∑}{\ensuremath{\textstyle\sum}}
\newunicodechar{∏}{\ensuremath{\textstyle\prod}}
\newunicodechar{∅}{\ensuremath{\emptyset}}
\newunicodechar{∪}{\ensuremath{\cup}}
\newunicodechar{∩}{\ensuremath{\cap}}
\newunicodechar{≡}{\ensuremath{\equiv}}
\newunicodechar{ℕ}{\ensuremath{\mathbb{N}}}
\newunicodechar{ℤ}{\ensuremath{\mathbb{Z}}}
\newunicodechar{ℝ}{\ensuremath{\mathbb{R}}}
\newunicodechar{𝔽}{\ensuremath{\mathbb{F}}}
\newunicodechar{α}{\ensuremath{\alpha}}
\newunicodechar{β}{\ensuremath{\beta}}
\newunicodechar{σ}{\ensuremath{\sigma}}
\newunicodechar{τ}{\ensuremath{\tau}}
\newunicodechar{φ}{\ensuremath{\varphi}}
\newunicodechar{ψ}{\ensuremath{\psi}}
\newunicodechar{χ}{\ensuremath{\chi}}
\newunicodechar{Φ}{\ensuremath{\Phi}}
\newunicodechar{Δ}{\ensuremath{\Delta}}
\newunicodechar{₀}{\textsubscript{0}}
\newunicodechar{₁}{\textsubscript{1}}
\newunicodechar{₂}{\textsubscript{2}}
\newunicodechar{ᵏ}{\ensuremath{{}^{k}}}
\newunicodechar{ⁿ}{\ensuremath{{}^{n}}}
\newunicodechar{ʲ}{\ensuremath{{}^{j}}}
\newunicodechar{²}{\ensuremath{{}^{2}}}
\newunicodechar{³}{\ensuremath{{}^{3}}}
\newunicodechar{⁴}{\ensuremath{{}^{4}}}
\newunicodechar{⁻}{\ensuremath{{}^{-}}}
\newunicodechar{¹}{\ensuremath{{}^{1}}}

% ---- Lean listing style ----
\definecolor{leanbg}{RGB}{247,247,250}
\definecolor{leankw}{RGB}{0,90,160}
\definecolor{leancm}{RGB}{120,120,120}
\lstdefinestyle{lean}{
  backgroundcolor=\color{leanbg},
  basicstyle=\ttfamily\small,
  keywordstyle=\color{leankw}\bfseries,
  commentstyle=\color{leancm}\itshape,
  morekeywords={theorem,lemma,def,fun,Prop,Type,variable,where,
    Field,Fintype,Finite,DecidableEq,CommRing,CommSemiring,Ring,CharP,Function,
    Injective,Surjective,Bijective,by,rfl,have,rw,exact,ext,simp,intro,
    iterateFrobenius,RingHom,Commute,convert,apply,using,Nat,Coprime,gcd,
    Odd,Even,IsAPN,IsAB,walsh,Nat,obtain,induction,cases,calc,grind},
  showstringspaces=false,
  columns=fullflexible,
  extendedchars=true,
  frame=single,
  framesep=6pt,
  xleftmargin=4pt,
  aboveskip=12pt,
  belowskip=14pt,
}
\lstset{style=lean}

\newtheorem{lemma}{Lemma}
\newtheorem{theorem}{Theorem}
\newtheorem{corollary}{Corollary}

% boxed "what just happened" note
\newcommand{\note}[1]{%
  \par\smallskip
  \noindent\colorbox{leanbg}{\parbox{\dimexpr\linewidth-2\fboxsep}{\small #1}}%
  \par\smallskip}


\title{\textbf{Two roots, and that is all}\\[4pt]
\large $x^2+x=0 \iff x\in\{0,1\}$ in characteristic $2$, read slowly}
\author{}
\date{}

\begin{document}
\maketitle

\begin{abstract}
\noindent
One line of algebra.\\[2pt]
$\bullet$\quad \texttt{sq\_add\_self\_eq\_zero\_char2}:\;\;
$u^2+u=0 \iff u=0 \lor u=1$.\\[6pt]
\noindent
It is the literal \emph{last step} of \texttt{kasami\_is\_apn}: once a collision
has been squeezed to $(x+y)^2+(x+y)=0$, this lemma delivers ``at most two
solutions'', i.e.\ APN.
\end{abstract}

\bigskip

% =====================================================================
\section*{Setting}

\[
\Gamma=\bigl(F:\mathrm{Type},\;[\mathrm{Field}\;F],\;[\mathrm{CharP}\;F\;2]\bigr)
\]
A field where $1+1=0$, so $-1=1$ and $u+u=0$ for every $u$.

\bigskip

% =====================================================================
\section{The lemma}

\[
\forall u:F.\qquad u^2+u=0 \;\longleftrightarrow\; (u=0 \;\lor\; u=1)
\]

\begin{lstlisting}
lemma sq_add_self_eq_zero_char2 {F : Type*} [Field F] [CharP F 2] {u : F} :
    u ^ 2 + u = 0 ↔ u = 0 ∨ u = 1 := by
  grind +suggestions
\end{lstlisting}

\note{\texttt{grind} closes it because the statement is a polynomial
identity over an integral domain: it factors and reads off the roots. We read
the factoring by hand below.}

\bigskip
\subsection*{Why it is true (the factoring)}

\noindent
Factor the left side:
\[
u^2+u \;=\; u(u+1).
\]
In a \textbf{field} (no zero divisors) a product is zero iff a factor is zero:
\[
u(u+1)=0 \;\Longleftrightarrow\; u=0 \;\lor\; u+1=0.
\]
In characteristic $2$, $\;u+1=0 \iff u=1\;$ (because $-1=1$). Hence
$u\in\{0,1\}=\mathbb{F}_2$.

\[
\begin{tikzcd}[column sep=large]
u^2+u=0 \arrow[r, leftrightarrow, "{\text{factor}}"]
& u(u+1)=0 \arrow[r, leftrightarrow, "{\text{domain}}"]
& u=0 \;\lor\; u=1
\end{tikzcd}
\]

\note{It is an Artin--Schreier polynomial $X^2-X = X^2+X$ in char $2$. Its
roots are exactly the prime subfield $\mathbb{F}_2$. The map $u\mapsto u^2+u$ is
$\mathbb{F}_2$-linear with kernel $\{0,1\}$ --- this is the additive shadow of
``$x\mapsto x^2$ is Frobenius''.}

\bigskip

% =====================================================================
\section{Where it lands: the last line of Kasami}

\noindent
Inside \texttt{kasami\_is\_apn} a collision between $x$ and $y$ is reduced,
via the trace identity, to
\[
x^2+x \;=\; y^2+y .
\]
Add the two sides and use Freshman's Dream $(x+y)^2=x^2+y^2$:
\[
(x+y)^2+(x+y) \;=\; (x^2+x)+(y^2+y) \;=\; 0 .
\]
Now the lemma fires on $u=x+y$:
\[
(x+y)^2+(x+y)=0
\;\xRightarrow{\;\texttt{sq\_add\_self\_eq\_zero\_char2}\;}\;
x+y=0 \;\lor\; x+y=1 .
\]
That is exactly $y=x$ or $y=x+1$ --- two solutions, which after rescaling by
$a$ is the APN condition.

\begin{lstlisting}
  have h_sum_zero : (x + y) ^ 2 + (x + y) = 0 := by
    have h2 : (x + y) ^ 2 = x ^ 2 + y ^ 2 := add_pow_char (R := F) (p := 2) x y
    rw [h2]
    have : x ^ 2 + y ^ 2 + (x + y) = (x ^ 2 + x) + (y ^ 2 + y) := by ring
    rw [this, hu, CharTwo.add_self_eq_zero]
  rw [sq_add_self_eq_zero_char2] at h_sum_zero
\end{lstlisting}

\note{Compare with the Frobenius pair. \texttt{frob\_bijective}+\texttt{frob\_add}
open the door (transport APN through a twist); this lemma closes it (read off the
two roots). Both are tiny; both are decisive.}

\bigskip

% =====================================================================
\section*{One-screen summary}

\begin{center}
\renewcommand{\arraystretch}{1.6}
\begin{tabular}{l l}
\textbf{step} & \textbf{content} \\
\hline
factor & $u^2+u=u(u+1)$ \\
domain & product $=0 \Rightarrow$ a factor $=0$ \\
char $2$ & $u+1=0 \iff u=1$ \\
land & $(x+y)^2+(x+y)=0 \Rightarrow y\in\{x,\,x+1\}$ \;(APN) \\
\end{tabular}
\end{center}

\bigskip
\noindent
\colorbox{leanbg}{\parbox{\dimexpr\linewidth-2\fboxsep}{\centering
$u(u+1)=0$ in a field of char $2$ $\;\Rightarrow\;$ Kasami collision has
$\le 2$ solutions $\;\Rightarrow\;$ APN.}}

\end{document}
